W niniejszym podrozdziale wprowadzono podstawowe oznaczenia wykorzystywane w pracy.

\begin{df}[Normy wektorowe i macierzowe \cite{golub1996matrix}]
Dla wektora $x \in \mathbb{R}^n$, jego normę $l_p$ dla $0 < p < \infty$ definiujemy jako 
$$\|x\|_{l_p^n} := \left(\sum_{i=1}^n |x_i|^p\right)^{1/p}.$$
Dla macierzy $X \in M_{n \times m}$, normę Frobeniusa definiujemy jako
$$\|X\|_F := \left(\sum_{i,j} |x_{ij}|^2\right)^{1/2}.$$
\end{df}

\begin{df}[Rozkład według wartości osobliwych (SVD) \cite{golub1996matrix}]
Dla dowolnej macierzy $X \in M_{n \times m}$ jej (zredukowany) rozkład SVD zapisujemy jako $$X = U \Sigma V^T,$$ gdzie $U \in M_{n \times p}$ oraz $V \in M_{m \times p}$ (przy $p \le \min(n,m)$) są macierzami o ortonormalnych kolumnach, a
$$\Sigma = \text{diag}(\sigma_1, \dots, \sigma_p),$$
jest macierzą diagonalną zawierającą wartości osobliwe $\sigma_1 \ge \sigma_2 \ge \dots \ge \sigma_p \ge 0$.
\end{df}

Zdefiniujmy dodatkowo klasę

$$ \mathcal{F}_d := \left\{ f: [0,1]^d \rightarrow \mathbb{R}, \lVert D^\alpha f\rVert_{\infty} \leq 1, \alpha \in \mathbb{N}^d_0\right\}$$

\subsection{Podstawy teorii Próbkowania Skompresowanego (Compressed Sensing)}

\begin{df}[Własność ograniczonej izometrii - RIP \cite{baraniuk2008simple}]
Własność ograniczonej izometrii (Restricted Isometry Property) określa, że dla macierzy $\Phi \in M_{m \times d}$ oraz pewnej stałej $0 < \delta < 1$, dla wszystkich wektorów $x \in \mathbb{R}^d$ o zadanej rzadkości (wektory wysokiego wymiaru, w których zdecydowana większość elementów ma wartość zero) zachodzi relacja: 
$$(1-\delta)\|x\|_{l_2^d}^2 \le \|\Phi x\|_{l_2^m}^2 \le (1+\delta)\|x\|_{l_2^d}^2.$$
\end{df}

\begin{df}[Najlepsza $K$-członowa aproksymacja \cite{wojtaszczyk2011}]
Dla wektora $x \in \mathbb{R}^d$, błąd najlepszej $K$-członowej aproksymacji w normie $l_1$ definiujemy jako $$\sigma_K(x)_{l_1^d} := \inf \{ \|x - z\|_{l_1^d} : \#\text{supp}(z) \le K \}.$$
\end{df}

\begin{tw}[Stabilna rekonstrukcja sygnału \cite{candes2006stable, wojtaszczyk2011}]
Załóżmy, że $\Phi$ jest losową macierzą wymiaru $m \times d$ (np. o elementach z rozkładu Bernoulliego). Przy odpowiednich założeniach co do wymiarów, rekonstrukcja $\hat{x} = \Delta(\Phi x + \epsilon)$ uzyskana za pomocą minimalizacji $l_1$ stabilnie przybliża wektor $x$, a błąd odzyskiwania $\|\hat{x} - x\|_{l_2^d}$ jest ograniczony przez sumę błędu najlepszej $K$-członowej aproksymacji $\sigma_K(x)_{l_1^d}$ oraz normy wektora szumu $\epsilon$ .
\end{tw}

\subsection{Geometria wysokich wymiarów i miary na sferze}

\begin{df}[Jednostajna miara na sferze \cite{ledoux2001concentration}]
Niech $\mu_{\mathbb{S}^{d-1}}$ oznacza jednostajną, rotacyjnie niezmienniczą miarę powierzchniową na sferze $\mathbb{S}^{d-1}$ w przestrzeni $\mathbb{R}^d$.
\end{df}

\begin{tw}[Rzut miary sferycznej na kulę nisko wymiarową \cite{rudin1980function}]
Niech $1 \le k < d$. Miara $\mu_k$ na jednostkowej kuli $B_{\mathbb{R}^k}$ w $\mathbb{R}^k$, będąca wynikiem rzutu miary $\mu_{\mathbb{S}^{d-1}}$ za pośrednictwem macierzy o ortonormalnych wierszach, wyraża się wzorem : 
$$d\mu_k(y) = \frac{\Gamma(d/2)}{\pi^{k/2}\Gamma((d-k)/2)}(1 - \|y\|_{l_2^k}^2)^{\frac{d-2-k}{2}}dy.$$
\end{tw}

\subsection{Kluczowe nierówności probabilistyczne}

\begin{tw}[Nierówność Hoeffdinga \cite{ledoux2001concentration}]
Niech $X_1, \dots, X_m$ będą niezależnymi zmiennymi losowymi, przyjmującymi prawie na pewno wartości z przedziałów $[a_j, b_j]$ . Wówczas dla sumy tych zmiennych prawdopodobieństwo odchylenia od wartości oczekiwanej o co najmniej $t$ jest ograniczone następująco : 
$$\mathbb{P}\left\{ \left| \sum_{j=1}^m X_j - \mathbb{E}\left(\sum_{j=1}^m X_j\right) \right| \ge t \right\} \le 2e^{-\frac{2t^2}{\sum_{j=1}^m (b_j - a_j)^2}}.$$
\end{tw}

\begin{tw}[Macierzowa nierówność Chernoffa \cite{tropp2012user}]
Rozważmy niezależne, losowe, dodatnio półokreślone macierze $X_1, \dots, X_m$ wymiaru $k \times k$, dla których maksymalna wartość osobliwa spełnia $\sigma_1(X_j) \le C$ prawie na pewno. Prawdopodobieństwo, że największa (lub najmniejsza) wartość osobliwa sumy $\sum_{j=1}^m X_j$ odchyli się o zadany czynnik od swojej wartości oczekiwanej, jest ograniczone wykładniczo. Twierdzenie to jest niezbędne przy szacowaniu zachowania sum macierzy losowych.
\end{tw}

\subsection{Stabilność podprzestrzeni niezmienniczych}

\begin{tw}[Twierdzenie Wedina o ograniczeniu perturbacji \cite{wedin1972perturbation}]
Niech dane będą macierze $B$ oraz $\hat{B}$ wraz z ich rozkładami SVD, gdzie $V_1$ i $\hat{V}_1$ reprezentują ich prawe wektory osobliwe odpowiadające największym wartościom osobliwym. Jeśli wartości osobliwe obu macierzy są odseparowane od siebie o co najmniej $\bar{\alpha} > 0$, to odległość podprzestrzeni rozpiętych przez wiersze można oszacować za pomocą normy Frobeniusa błędu: 
$$\|V_1 V_1^T - \hat{V}_1 \hat{V}_1^T\|_F \le \frac{2}{\bar{\alpha}}\|B - \hat{B}\|_F.$$
\end{tw}